

质子交换膜燃料电池(PEMFC)作为高效清洁的能源转换装置，在新能源汽车和分布式发电领域具有广阔的应用前景。气体扩散层(GDL)作为其核心组件，由碳纤维和粘结剂构成，其微观结构直接决定了气体传输、水管理和电池整体性能。因此，精确重构GDL的三维微观结构，对于理解其内部传质机理、优化材料设计具有重要的理论与现实意义。

目前，国内外学者已提出多种GDL微观结构重构方法，主要包括基于实验图像的重构和基于随机算法的生成式重构。然而，现有研究在GDL三维重构中仍存在显著局限：多数随机重构模型不仅将碳纤维简化为理想的长直圆柱体，还进一步假设其在整个厚度方向上均匀分布，完全忽略了实际生产中纤维的交错重叠与局部密度不均(Zhang et al, 2022；Gao et al, 2020)。这种双重简化导致重构的孔隙结构与真实GDL存在明显偏差，进而影响了传质模拟的准确性。同时，X射线CT等实验技术也难以有效区分碳纤维与粘结剂相(Benz et al, 2025；Jin et al, 2019)，导致对其分布的定量分析存在困难。由于上述简化，现有模型往往无法准确预测GDL在压缩等工况下的内部结构变化(Wang et al, 2022)，使得模拟得到的气体渗透、水管理等性能与实验结果存在偏差，限制了模型对GDL设计优化的指导能力。

尽管已有研究尝试通过分层堆叠的方式模拟GDL的厚度方向结构（Lai et al, 2026），但大多仅调控各层的孔隙率，而未对纤维的空间取向与位置分布进行精细化控制。为了在计算可行性与物理真实性之间取得平衡，本研究在现有随机重构模型的基础上，重点突破“均匀分布”的假设，提出了分层取向与分层位置采样策略，对纤维的空间分布与取向进行更精细的建模，并进一步引入粘结剂的分布模型，以更贴近真实GDL的复杂结构。

本研究的目标是：(1)重构出给定孔隙率的碳纤维（重叠/不重叠）气体扩散层三维结构；(2)考虑粘结剂(binder)的影响，重构出更贴近真实情况的GDL三维结构。通过本研究，旨在提供一套完整的、可执行的三维重构流程与代码，并生成相应的三维数据文件，为后续的传质模拟和性能分析提供可靠的基础模型。

